\chapter{Estado del arte:}

Para el desarrollo de la sección del estado del arte se han trabajado en los siguientes campos, investigación de desarrollo web, gestión de usuarios y credenciales, entornos de virtualización y proceso de trazabilidad de firmas.

\section{Desarrollo de Páginas Web}

El campo de desarrollo de las páginas web han estado en auge desde su llegada en los 90, con la llegada de \texttt{HyperText Markup Language}, donde en un principio solo permitía la marcación básica del contenido. En este punto, las páginas web eran de contenido estático hasta que aparecieron los CSS que permitieron estilizar las páginas web destacando lenguajes como \texttt{Javascript} y \texttt{Personal Home Page}. En la década del 2010 aparecieron diferentes frameworks tanto para frontend como para Backend donde se puede destacar herramientas como AngularJS y React.js.\cite{armadillo_amarillo_evolucion_web} 


Hoy en día existe una gran multitud de posibilidades para el desarrollo web,
desde lenguajes como \texttt{TypeScript}, \texttt{JavaScript} o \texttt{Python}
hasta frameworks especializados que abstraen la complejidad del desarrollo de
interfaces. En este contexto, la elección del framework de frontend es una
decisión importante, ya que condiciona la mantenibilidad, la escalabilidad y la
experiencia de desarrollo a lo largo de todo el proyecto.

Entre las opciones más consolidadas destacan \texttt{Angular}, \texttt{Vue.js}
y \texttt{React.js}. Angular es un framework completo desarrollado por Google
que impone una estructura rígida, adecuada para equipos grandes pero con una
curva de aprendizaje elevada. Vue.js es una alternativa más ligera y progresiva,
pero con un ecosistema menor. React.js, desarrollado por Meta en 2013,
se ha consolidado como la biblioteca de frontend más utilizada a nivel mundial
gracias a su modelo basado en componentes reutilizables y su
Virtual DOM, que minimiza las actualizaciones reales del navegador
optimizando el rendimiento de la interfaz.

En el contexto de este proyecto, React.js aborda diferentes necesidades.
En primer lugar, el panel de administración de este proyeccto, requiere múltiples vistas
con estados complejos e interdependientes, como la gestión de usuarios, la
monitorización de colas de tareas o el seguimiento de solicitudes de acceso;
el modelo de componentes de React permite encapsular cada una de estas vistas
de forma independiente y reutilizarla en distintos contextos. En segundo lugar,
el uso de \texttt{TypeScript} sobre React añade tipado estático, lo que reduce
errores en tiempo de desarrollo y facilita el mantenimiento del código a lo
largo del tiempo.


% Tras una introducción histórica, vamos a introducir el estándares de industria para cada uno de los campos. En el ámbito de la programación de páginas web se distingen las siguientes tecnologías distinguiendo para el backend, bases de datos y frontend. Para el backend, se plantean diferentes tecnologías como NodeJs, Python, PHP y Java entre otras; en el caso de este proyecto, se utiliza NodeJs con diferentes librerías. Para bases de datos existen las relacionales y las no relacionales, aprovechando cada sistema se usan para cada tarea. Por ejemplo, en el caso de datos relacionados y estructurados se ha usado una relacional, en cambio, para una cola de mensajes donde los datos no son estructurados y se necesita una cola de mensajes para un paso ágil de los mismos se hace uso de una base de datos no relacional. En el frontend existen diferentes tecnologías que se pueden combinar, en este caso, se hace uso de Next.js junto a CSS entre otras. \cite{proun_tecnologias_web} 


\section{Gestión de usuarios y credenciales:}

En el ámbito de gestión de usuarios y credenciales, de igual manera han tenido una gran evolución, en la decada de los 60, el sistema predominante de gestión de usuario era un sistema de credenciales básicos, es decir, un usuario y una contraseña. En los 90, las redes empresariales empezaron a expandirse y se enfrentaron a nuevos retos de seguridad y apareció el sistema de \texttt{Self-service Password Reset} permitiendo el sistema de recuperación de password. A principios de los 2000 se introduce el \texttt{Single Sign-On (SSO)} incorporando un sistea de inicio único de sesión y herramientas de \texttt{Identity-as-a-Service} para gestionar sistemas distribuidos desde la nube. En el 2010 entra en escena la moda de \texttt{Bring Your Own Device} (BYOD) con nuevos riesgos de seguridad, para mitigarlas aparecen tecnologías como herramientas multifactor, token y accesos biométricos. En el caso de este sistema permite la conexión desde cualquier dispositivo mediante el uso de credenciales válidas. \cite{optimalidm_iam_evolution}

% Para la gestión de usuarios y credenciales, con el objetivo de tener una fuente única de verdad, en la base de datos se incorporan las credenciales usando en el sistema de login \texttt{JSON Web Token} y un sistema de altas por invitación. 


En las soluciones autoalojadas Keycloak es la referencia más extendida. Este sistema ofrece SSO, federación de identidades e inicio de sesión social,
y su naturaleza de código abierto permite una personalización extensiva y un
desarrollo impulsado por la comunidad.  No obstante,
el coste total de operación de Keycloak es elevado debido
al trabajo operacional, el mantenimiento y la necesidad de talento
especializado. En
el contexto de este trabajo para un laboratorio con alta rotación de personal y sin
un equipo DevOps dedicado, asumir el mantenimiento de una plataforma no es viable. Además el sistema no solo
se necesita autenticar usuarios, sino gestionar su ciclo de vida
completo, asociarlos a proyectos, aprovisionar automáticamente los servicios
que necesitan y revocarlos al desvincularse. Este TFG adopta un enfoque propio
sobre PostgreSQL como fuente única de verdad, utilizando
JSON Web Tokens (JWT) para la autenticación sin estado en la API y
un sistema de altas por invitación que garantiza que únicamente el personal
autorizado puede registrarse en la plataforma. Esta solución a medida elimina
la sobrecarga operativa de las alternativas analizadas y se adapta exactamente
a las necesidades del laboratorio.


\section{Virtualización}

El incio de la virtualización se remonta a los 60 donde se desarrollaron equipos de cómputos muy potentes como los 'mainframes', que debido a su alta capacidad, cierta parte del cómputo quedó sin utilización, entonces IBM desarrolló un método para ejecutar procesos en paralelo naciendo así las máquinas virtuales. No fue hasta 1999 hasta que se dió el siguiente paso, cuando \texttt{VMware} presentó su primer producto, donde se podían ejecutar varios sistemas operativos en la misma máquina. En 2001 nace los hipervisores de tipo 1, es decir, aquellos que son capaces de ejecutar directamente sobre el hardware teniendo acceso al mismo. Durante los siguientes años se desarrollaron diversas herramientas para la virtualización, como la aparición de VirtualBox, Microsoft-Hyper-V de máquinas, pero, el siguiente gran hito no fue hasta la aparición de AWS en el año 2012  haciendo uso de la virtualización. \cite{icorp_virtualizacion_ti} 


% Actualmente para la gestión de microservicios on premises las principales opciones son \texttt{Docker} o \texttt{Kubernetes}. Docker destaca por su capacidad de crear imágenes para realizar tareas complejas y relacionarlos con composiciones, mientras que Kubernetes es para orquestar clústers de microservicios. En el ámbito de esta práctica no es necesario trabajar con clústers, por lo que se decide trabajar únicamente con docker. A lo largo de este proyecto se van a usar diferentes tipos de imágenes para el despliegue de servicios, destacando entre otros Postgres y Redis para las bases de datos, Node para frontend y backend; y por último, code-server para entornos virtualizados de desarrollo. \cite{aws_kubernetes_vs_docker}

En el contexto de este proyecto, la virtualización responde a dos necesidades
concretas e independientes.

La primera es el aislamiento de los entornos de desarrollo de los
investigadores. Uno de los problemas identificados en el centro era que la
instalación de librerías y controladores por parte de distintos investigadores
sobre los servidores compartidos generaba conflictos de dependencias y
corrupciones del sistema operativo. La solución adoptada es proporcionar a
cada miembro del personal docente e investigador un contenedor
\texttt{code-server} propio, un entorno de desarrollo accesible desde el
navegador completamente aislado del resto. De este modo, las instalaciones
de un investigador no pueden interferir con las de otro ni comprometer la
estabilidad del servidor anfitrión.

La segunda es el despliegue reproducible de la propia plataforma.
Todos los servicios que componen la aplicación web, incluyendo las bases de
datos \texttt{PostgreSQL} y \texttt{Redis}, el backend en \texttt{Node.js}
y el frontend, se despliegan como contenedores Docker orquestados mediante
\texttt{docker-compose}. Esto garantiza que el entorno de producción es
reproducible en cualquier máquina, elimina la dependencia de configuraciones
manuales del sistema operativo anfitrión y simplifica tanto el despliegue
inicial como el mantenimiento posterior.

Se descarta \texttt{Kubernetes} por sobredimensionamiento: está diseñado
para escenarios con múltiples nodos, alta disponibilidad y escalado
automático, requisitos que exceden el alcance de un laboratorio de
investigación de tamaño reducido. \texttt{Docker} con
\texttt{docker-compose} cubre ambas necesidades del proyecto con una
complejidad operativa significativamente menor y sin requerir infraestructura
adicional.


\section{Firma electrónica}


En España en 1955, nació la primera Autoridad de Certificación y Europa en 1999 estableció la primera normativa de las firmas digitales.En el 2003 nace la primera legislación Española de la firma electrónica. Tres años después, en el 2006, nace el DNIe permitiendo a la ciudadanía identificarse de manera electrónica y realizar operaciones y consultas. Posteriormente en el 2014 se instaura el reglamento eIDAS en toda la unión europea, un marco legal de identificación electrónica y servicios de confianza. Además en 2015, nace el DNI 3.0 permitiendo su lectura con NFC. Por último en 2024 entra en vigor eIDAS 2, introduciendo un monedero digital que permite guardar las credenciales digitales junto a otros datos personales. \cite{viafirma_evolucion_firma}

En el contexto de este proyecto, la firma electrónica se emplea como mecanismo
de trazabilidad y validez legal para las solicitudes de acceso a los servidores
y servicios del laboratorio. El requisito fundamental es que dichas firmas sean
jurídicamente válidas en España, lo que implica reconocer los certificados
emitidos por la Fábrica Nacional de Moneda y Timbre (FNMT) y los
asociados al DNIe, emitidos por la Dirección General de la Policía Nacional.
La firma digital avanzada o cualificada, basada en un
certificado digital como el de la FNMT o el DNIe, es la única con validez
para trámites oficiales. 

La primera opción considerada fue AutoFirma, la herramienta
oficial del Ministerio de Transformación Digital y de la Función Pública para la firma de documentos en trámites
con la Administración. Si bien reconoce los certificados de la FNMT y el DNIe,
está diseñada como aplicación de escritorio para uso ciudadano, no como
biblioteca integrable en una plataforma web, lo que la descarta para este caso
de uso. La solución adoptada es el Digital Signature Service (DSS), una
biblioteca de código abierto desarrollada y mantenida por la Comisión Europea.
DSS soporta múltiples formatos de firma, incluyendo PAdES,
XAdES entre otras. Su ventaja fundamental frente a las alternativas es que
incorpora gestión de confianza mediante la carga de listas
de confianza europeas.  Esto significa que la plataforma puede validar
automáticamente los certificados de la FNMT y el DNIe sin necesidad de
mantener manualmente las listas de autoridades de certificación reconocidas,
eliminando el riesgo de desactualización y garantizando el cumplimiento
continuo con la normativa eIDAS. \cite{eu_digital_signature_service}






